\subsection{D1 interim: Jig's defences on AgentDojo}
\label{sec:results-d1}

Four of the five planned configurations have finished all 140 of their trials (560 of 700); the fifth,
the Sentinel on the Q4\_K\_M Qwen counterpart, could not start because that model needs about
21.7\,GB with the safety margin beside the owner's server and 18.2\,GB was free, so it remains queued.
No finished trial is a harness error (14 errors in the first attempt, caused by a pipeline-naming fault,
were retried; see the deviations log). Table~\ref{tab:d1-d1-full-v1} gives every cell. The run has one seed
and is incomplete, so we report intervals and defer hypothesis tests to the complete multi-seed run, as
the protocol requires.

\paragraph{Attack success is at the floor without any defence.} Pooled over both attacks, the attack
succeeded in 3/120 injected trials with no reviewer (2\%; Wilson 95\% interval 1--7\%), 4/120 with the
Sentinel on the agent's own model (3\%; 1--8\%), 4/120 with Granite Guardian (3\%; 1--8\%) and 2/120 in
read-only mode (2\%; 0--6\%). Every one of these 13 successes is the same AgentDojo case, travel
\texttt{injection\_task\_6}, whose goal is met by what the agent \emph{says} (its ground truth needs no tool
call). On the 110 injected trials per configuration whose goal needs a tool call, the attack never
succeeded in any configuration (0/110; 0--3\%). With this agent (the ternary Bonsai 2 27B at temperature 0)
and this subsample, the undefended agent almost never acts on an injection, so the benchmark cannot
separate the defences on security: H-D1 and H-D2 are neither supported nor refuted here.

\paragraph{Defences differ in what they cost.} With no attack present, the user's task succeeded in
16/20 trials with no reviewer (80\%; 58--92\%), 14/20 with the Sentinel (70\%; 48--85\%), 11/20 with
Granite Guardian (55\%; 34--74\%) and 8/20 in read-only mode (40\%; 22--61\%). Under attack, on six times as
many trials, the Sentinel matched the no-reviewer level while Guardian and read-only mode again lost
utility: 92/120 (77\%; 68--83\%) with no reviewer, 93/120 (78\%; 69--84\%) with the Sentinel, 71/120
(59\%; 50--68\%) with Guardian and 55/120 (46\%; 37--55\%) read-only. Read-only mode's loss is expected (it refused 69 side-effecting
calls that tasks needed), and is the cost side of H-D3; Guardian's loss of about 18 points under attack
bought no measurable security on this subsample. The Sentinel on the agent's own model kept utility at the
no-reviewer level but escalated 42 calls to the approval queue over 140 trials (0.3 per task), all of
which the labelled \texttt{approve\_all} owner approved; Guardian, which only allows or denies, raised none.

\paragraph{Most blocks come from a core rule, not a reviewer.} Recording which layer refused each call
(available for trials after 04:55; all 140 trials for both Sentinel configurations, 87 of 140 with no
reviewer, none for Guardian) shows that Jig's core rule \texttt{no-local-network} refused far more calls
than any reviewer: 239 refusals against 6 Sentinel denials in action mode, and 176 against 10 in read-only
mode. The rule rejects URLs without an \texttt{http}/\texttt{https} scheme, and AgentDojo's Slack suite writes
web addresses as bare host names (\texttt{www.\ldots}), so the rule blocks benign and injected web calls
alike. It is a genuine Jig behaviour, but an incidental one: it lowers Slack utility in every
configuration and would also stop any injection whose goal is to fetch or post to such an address.
